¿Cómo se puede diseñar e implementar un sistema de apoyo a la decisión que, mediante modelos de predicción de demanda, permita estimar las cantidades necesarias de productos en un sistema de gestion de inventarios?

% \subsubsection{ESTRUCTURACION DE LA PREGUNTA MEDIANTE PICO}

% Con el fin de estructurar de manera formal la pregunta de investigación, se empleo la metodología PICO (Population, Intervention, Context, Outcome). Esta metodología ayuda a  identificar los elementos clave del problema y orientar el desarrollo de la solución propuesta.

% \begin{itemize}
%     \item \textbf{P (Population):} Sistemas de gestión de inventario de diferentes dominios.
    
%     \item \textbf{I (Intervention):} Implementación de un sistema de apoyo a la decisión basado en modelos de predicción de demanda.
    
%     \item \textbf{C (Context):} Proceso de planificación y adquisición de productos dentro de un sistema de gestion de inventarios.
    
%     \item \textbf{O (Outcome):} Generación de estimaciones de demanda que apoyen la toma de decisiones en la gestión de inventarios.
% \end{itemize}

\subsubsection{Preguntas Derivadas}

¿Qué variables y datos históricos son necesarios para estimar la demanda mensual de productos?

¿Qué modelos de predicción de demanda son más adecuados para estimar el consumo o venta de productos en diferentes dominios?

¿Cómo debe estructurarse el proceso de análisis y procesamiento de datos para generar predicciones confiables?

¿Qué tipo de salidas debe generar el sistema (por ejemplo, cantidades estimadas, punto de reposición, alertas) para apoyar la toma de decisiones?
